\chapter{INTRODUCCIÓN}

\introsection{Presentación del Proyecto}

El empleo de cinemática bioinspirada en robots caminantes busca sortear las limitaciones de tracción que presentan las plataformas convencionales de ruedas u orugas sobre terrenos no estructurados~\cite{James2021_InsectLegDesign, DungBeetle_2023_Cybernetics}. En la minería subterránea, este desafío mecánico coincide con un entorno de alto riesgo ocupacional: las galerías estrechas concentran emanaciones de monóxido de carbono ($\text{CO}$) y óxidos de nitrógeno ($\text{NO}_x$) generados por la combustión de maquinaria pesada, donde una ventilación deficiente eleva la probabilidad de intoxicación del personal o paraliza las operaciones~\cite{AdvancesMining_2023, AdaptiveControl_Mines_2025}. En el contexto peruano, la heterogeneidad entre la minería formal y la de pequeña escala acentúa la falta de protocolos de monitoreo automatizados y accesibles. Frente a ello, los esquemas robóticos con morfología de insecto permiten distribuir puntos de apoyo en sustratos rocosos inestables y actuar como soporte instrumental para tareas de inspección y registro ambiental, evitando la exposición directa de cuadrillas a atmósferas no respirables~\cite{KimChoi_2025_MineMapping, Zhou2023_BioinspiredSensors, Liu2025_PiezoBioinspired}.


\introsection{Descripción de la situación problemática}

En las labores mineras subterráneas, las actividades de perforación, voladura y acarreo de mineral se realizan principalmente con maquinaria pesada diésel, tales como jumbos perforadores, scooptrams y volquetes de bajo perfil. Sus gases de combustión, principalmente monóxido de carbono ($\text{CO}$), dióxido de carbono ($\text{CO}_2$) y óxidos de nitrógeno ($\text{NO}_x$), alteran rápidamente la composición del aire en frentes de trabajo con ventilación restringida~\cite{Rodriguez2023_EconomiaCircularMinera, TorresSalazar2021_Diesel}. En las operaciones mineras del Perú, estas emisiones demandan caudales de aire forzado continuos para evitar concentraciones letales~\cite{Cortez2016_GEI_Minas}. Conforme el avance de las labores profundiza en nuevas galerías o cruceros ciegos, el circuito de ventilación pierde eficiencia, lo que acelera la acumulación de contaminantes gaseosos y reduce el porcentaje volumétrico de oxígeno disponible. Desde la perspectiva de la ingeniería mecatrónica, el problema combina restricciones de flujo aerodinámico, geometrías de sección reducida y severas limitaciones para el despliegue de instrumentación fija en espacios confinados. Por estas razones, tanto los informes técnicos de supervisión como el marco normativo minero nacional exigen garantizar condiciones ambientales seguras antes del ingreso de cualquier cuadrilla~\cite{OSINERGMIN_BoletinGSM_2020, MINEM_DS024_2016}.

La exposición prolongada a estas atmósferas genera intoxicaciones agudas, hipoxia y daños pulmonares severos si no se detectan a tiempo. A nivel internacional y local, se reconoce que la cuantificación precisa de concentraciones bajas de $\text{CO}$ y $\text{NO}_2$ en galerías subterráneas continúa siendo un desafío metrológico debido al polvo en suspensión y a la humedad relativa~\cite{EU_ExposureDirective_2020}. Por ley, toda operación requiere mediciones atmosféricas previas al inicio de jornada; sin embargo, las auditorías técnicas reportan inspecciones ejecutadas de forma tardía o sin la rigurosidad debida, derivando en fatalidades evitables durante eventos de voladura o fallas de tiro~\cite{MINEM_Yanaquihua_2023, OSINERGMIN_BoletinGSM_2020, MINEM_DS024_2016}.


\introsection{Formulación del problema}

En las labores mineras subterráneas peruanas, el personal técnico se expone de forma recurrente a ambientes contaminados debido a una ventilación auxiliar insuficiente y a la acumulación de gases tóxicos. Dicha situación, documentada en los informes de supervisión de OSINERGMIN y regulada por el Decreto Supremo N$^\circ$ 024-2016-EM del Ministerio de Energía y Minas (MINEM), constituye un foco permanente de riesgo laboral~\cite{OSINERGMIN_BoletinGSM_2020, MINEM_DS024_2016}. Pese a la fiscalización vigente, la siniestralidad reciente demuestra que la inhalación de atmósferas no respirables sigue causando accidentes graves en minas de socavón~\cite{MINEM_Yanaquihua_2023}.

En la práctica operativa, el monitoreo del aire recae en cuadrillas que recorren las galerías con detectores portátiles multigás. Si bien el protocolo cumple formalmente la regulación, obliga a los inspectores a ingresar físicamente a frentes ciegos antes de certificar si la atmósfera es respirable, exponiéndolos a concentraciones peligrosas si el tiro de ventilación ha fallado o si los gases de voladura no se han disipado~\cite{MINEM_DS024_2016, OSINERGMIN_BoletinGSM_2020}. Asimismo, la precisión de los sensores electroquímicos portátiles se degrada en presencia de humedad extrema y material particulado, lo que dificulta discernir a tiempo niveles traza de $\text{CO}$ y $\text{NO}_2$ críticos para la seguridad del personal~\cite{EU_ExposureDirective_2020}.

Frente a estas restricciones operativas, se requiere transferir la tarea de muestreo ambiental a plataformas autónomas o teleoperadas capaces de ingresar de forma remota durante las inspecciones preliminares. Reemplazar la presencia humana en la primera línea de medición permite verificar el estado de los frentes de trabajo y registrar variables ambientales sin comprometer la integridad física de las brigadas~\cite{OSINERGMIN_BoletinGSM_2020, MINEM_DS024_2016}.

En este contexto, la inspección remota en labores subterráneas puede resolverse mediante diversas alternativas tecnológicas, incluyendo vehículos terrestres convencionales, sistemas aéreos no tripulados o arquitecturas robóticas cooperativas. En el Capítulo II se evalúan cuantitativamente estas alternativas frente a criterios de estabilidad mecánica, autonomía energética y restricciones ambientales propias del socavón minero.


\introsection{Objetivos del Proyecto}

\textbf{Objetivo general} \\
Diseñar un sistema robótico cooperativo bioinspirado en insectos, conformado por dos unidades autónomas, capaz de realizar mapeo del terreno y detección de gases tóxicos en zonas de difícil acceso en minería subterránea, integrando criterios derivados de la literatura reciente y de la normativa vigente en ventilación, seguridad ocupacional y monitoreo ambiental.

\textbf{Objetivos específicos}
\begin{itemize}
    \item Desarrollar e implementar la arquitectura mecánica del sistema robótico, mediante la adaptación de la unidad principal y el diseño de la unidad secundaria garantizando estabilidad operacional y resistencia ante perturbaciones externas.
    \item Diseñar e implementar la placa de control basada en un chip electrónico para los motores.
    \item Desarrollar los algoritmos de control, navegación y comunicación que permitan la coordinación cooperativa entre las dos unidades durante las tareas de exploración.
    \item Desarrollar el sistema de monitoreo y detección de gases tóxicos.
    \item Integrar y evaluar experimentalmente el sistema cooperativo en condiciones controladas de laboratorio, verificando la efectividad en la detección de gases tóxicos, la precisión del mapeo y la estabilidad de la comunicación entre unidades.
\end{itemize}

La justificación de este proyecto responde a la recurrencia de incidentes por asfixia e intoxicación en labores subterráneas peruanas, donde los frentes ciegos carecen de ventilación continua~\cite{OSINERGMIN_BoletinGSM_2020, MINEM_Yanaquihua_2023}. En lugar de enviar inspectores a zonas no verificadas, un sistema mecatrónico autónomo ejecuta la exploración inicial y transmite mapas de concentración de gases. La adopción de una morfología bioinspirada en insectos responde a un criterio mecánico concreto: las extremidades multiarticuladas proporcionan estabilidad estática y facilitan el franqueo de roca suelta, zanjas y pendientes pronunciadas donde los vehículos convencionales patinan o vuelcan~\cite{InsectInspired_2020, DungBeetle_2023_Cybernetics}.

Tanto la literatura especializada en emisiones diésel como los reportes de calidad del aire subterráneo confirman que la dilución de gases tóxicos en frentes de avance es irregular y propensa a acumulaciones críticas~\cite{TorresSalazar2021_Diesel, Cortez2016_GEI_Minas, Cepa2023_sensors}. A su vez, el Reglamento de Seguridad y Salud Ocupacional en Minería (D.S. 024-2016-EM) y las directivas de fiscalización de OSINERGMIN exigen el registro auditable y la verificación previa de atmósferas antes de autorizar cualquier labor~\cite{OSINERGMIN_BoletinGSM_Ventilacion_2020, MINEM_DS024_2016, MINEM_Yanaquihua_2023}. Por tanto, la automatización del sensado no constituye únicamente una ventaja de rapidez operativa, sino un mecanismo para garantizar la trazabilidad de las mediciones conforme a los límites de exposición ocupacional.

Para maximizar la cobertura sin comprometer la maniobrabilidad, el sistema se estructura bajo una configuración cooperativa integrada por dos plataformas con funciones desacopladas. Esta estrategia emula la división funcional observada en enjambres biológicos~\cite{InsectInspired_2020}: una unidad de mayor escala transporta el equipamiento de cómputo, cartografía 3D y enlace inalámbrico, mientras que una unidad secundaria de menor sección transversal se desprende para explorar cavidades angostas o con derrumbes parciales~\cite{FlappingFlight_2021}. Esta arquitectura distribuida extiende el radio de inspección sin sobrecargar cinemáticamente a un único robot, garantizando mediciones georreferenciadas y continuas en puntos donde el personal no puede ingresar.


\introsection{Alcance y limitaciones}

Este proyecto comprende el diseño constructivo, la manufactura y la validación experimental de dos hexápodos cooperativos en banco de pruebas y pistas de laboratorio que simulan pendientes e irregularidades mineras. El desarrollo abarca el dimensionamiento cinemático de las extremidades, el diseño de la electrónica modular para el control articular y de potencia, el enlace de comunicación local entre ambas unidades y la integración del paquete de sensado para gases ($\text{CO}$ y combustibles). La validación se centra en corroborar la estabilidad de la marcha en trípode, la repetibilidad en la transmisión de telemetría y la coherencia del mapeo espacial bajo condiciones controladas.

Por otro lado, el desempeño de la cooperación entre las unidades robóticas estará estrictamente condicionado por la disponibilidad, alcance, latencia y estabilidad del enlace de comunicación inalámbrico. La autonomía operativa del sistema estará limitada por la capacidad energética de las baterías frente al consumo acumulado de los sistemas de actuación, sensado, cómputo y comunicación. Finalmente, la investigación se limita al desarrollo de un prototipo a nivel experimental, por lo que el sistema no está diseñado ni certificado normativamente para su operación inmediata en ambientes clasificados como peligrosos o minas subterráneas reales~\cite{OSINERGMIN_DS024_Compilado, OSINERGMIN_BoletinGSM_2020}.
